% =====================================================================
%  MODELLO: LETTERA  (lettera di presentazione per RTB/PB, lettera di
%  candidatura, comunicazioni ufficiali)
%
%  1. Una lettera è un documento esterno: copiate questo file in
%       <Milestone>/documenti/esterni/   es.
%       Candidatura/documenti/esterni/lettera-candidatura.tex
%       RTB/documenti/esterni/lettera-presentazione-rtb.tex
%  2. Cambiate i campi qui sotto (dopo il simbolo % c'è un commento).
%  3. Scrivete il testo dopo \begin{document}.
%  Carta intestata, data, destinatari, oggetto, saluti e firma si fanno
%  da soli. La lettera non ha registro delle modifiche: la verifica si
%  fa con la pull request. Guida completa: GUIDA.md (ricetta 6).
% =====================================================================
\documentclass[lettera]{gpt}

\titolo{Lettera di presentazione}  % nome della lettera (sito e PDF)
\oggetto{Oggetto della lettera}    % riga "Oggetto:" (se manca, usa il titolo)
\data{2026-10-26}                  % data della lettera: AAAA-MM-GG
\luogo{Padova}                     % città accanto alla data
\responsabile{spagnoli}            % chi firma: il Responsabile del gruppo

% Destinatari: se non scrivete niente valgono quelli di dati-gruppo.tex
% (i due docenti). Per scrivere ad altri, togliete il % qui sotto:
% \destinatari{Spett.le Azienda S.p.A.\newline Via Roma 1\newline 35100 Padova}

% Saluto finale (facoltativo, questo è quello predefinito)
% \chiusura{Cordiali saluti,}

\begin{document}

Egregio Prof.~Vardanega, egregio Prof.~Cardin,

con la presente il gruppo \textbf{Git Push Team} vi comunica ...

Il materiale è disponibile al seguente indirizzo:
\url{https://git-push-team.github.io/docs/}

\end{document}
